%=====================================================================
%  SECTION: What the jobs model buys that the continuous model cannot
%  Drop-in for "Jobs and Well-Being Measurement" (Haydar & Maniquet 2025)
%  Notation follows the draft: bundles z=(c,j), j in J, ability set
%  A \subseteq J, pre-tax income profile y: J -> R_+, preferences R over
%  (c,j), well-being measure W(z,R,A;y). Measures W^3 (own-A laisser-faire)
%  and W^{5,\bar A} (reference-ability laisser-faire).
%
%  Every substantive claim is annotated with its support, traceable to
%  INDEX_Motivation_Taxation_SWF.md. Three cautions from that index are
%  built in: (A1) we never say the continuous model "cannot carry
%  attributes"; (A4) genuinely new features are flagged as new; (B1/B3/B4)
%  aggregation and tax shapes are presented as precedented/natural, NOT as
%  theorems of this model.
%=====================================================================

\section{What the model adds to the classical labour--leisure benchmark}
\label{sec:contribution}

The standard normative model of labour income taxation, descending from
\textcite{mirrlees1971}, represents an individual's work by a scalar labour
supply $\ell\in[0,1]$ and an exogenous wage $w_i$, so that pre-tax income is
$w_i\ell$ and the only behavioural margin is \emph{how much} the individual
works. The fair-allocation refinement of this model
\textcite{FleurbaeyManiquet2006,FleurbaeyManiquet2007,fleurbaey_maniquet_2018}
retains the same primitive: heterogeneity is carried by the single index
$w_i$, and a worker's situation is a point on a one-dimensional, fully
divisible, completely ordered labour line.

Our model replaces that line with a discrete set of jobs. An individual is
described by preferences $R$ over consumption--job bundles $(c,j)$, an
ability set $A\subseteq\mathcal J$ of jobs she can hold, and a pre-tax income
profile $\mathbf y:\mathcal J\to\mathbb R_+$ that prices each job. This is not
a cosmetic change of variables. We first state precisely what the
substitution does and does not claim against the classical model
(\S\ref{sub:claim}), then draw out the four consequences that motivate the
exercise: the separation of opportunity from choice
(\S\ref{sub:separation}), the relocation of the compensation--responsibility
cut (\S\ref{sub:compresp}), the appearance of feasibility phenomena with no
continuous analogue (\S\ref{sub:feasibility}), and the implications for how a
well-being measure should enter social evaluation and taxation
(\S\ref{sub:swf}).

\subsection{What is and is not being claimed}
\label{sub:claim}

It would be wrong to claim that the continuous model cannot accommodate
non-pecuniary job characteristics; one can always append attributes to the
pair $(c,\ell)$. The claim is narrower and structural. The continuous model
does not represent a \emph{heterogeneous feasible set of jobs} that differs
across individuals, and it does not distinguish \emph{which} categorical,
unordered job an individual holds---only the quantity of labour she supplies.
Two jobs with identical hours and identical pay but different content are, in
the labour-line representation, the same point.\footnote{%
\textcite{fleurbaey2008income}, \S5.7, explicitly identifies job heterogeneity
as a limitation of the consumption--leisure model and offers informal
remedies (absorbing job unpleasantness into a zero-equivalent reference, or
assuming the unskilled share similar jobs);
\citet{fleurbaey_maniquet_2019} confront the issue directly by treating
employment and labour as discrete, non-classical goods that violate both
monotonicity and cardinal divisibility. Our contribution is to make the job a
primitive carrying fixed non-pecuniary attributes, rather than to patch the
labour-line model.}
That a job is categorical---neither ``more is better'' nor convexly
combinable with another job---is exactly the pair of properties
(\citealp{fleurbaey_maniquet_2019}, on relaxing monotonicity and cardinality)
that the labour scalar $\ell$ is built to satisfy and that the job index $j$
is built to deny. We model preferences over $(c,j)$ so that
$(c,j)\mathrel{I}(c,j')$ may fail even when $\mathbf y(j)=\mathbf y(j')$:
equal pay need not mean equal well-being, because the jobs themselves differ.

\subsection{Opportunity and choice become separate primitives}
\label{sub:separation}

The central structural gain is the disentangling of three objects that the
classical wage $w_i$ fuses into one. In the Mirrleesian model the wage
simultaneously encodes what the individual \emph{can} do (her productive
opportunity) and indexes the budget line along which she \emph{chooses}
(her position $w_i\ell$); pay and productivity are the same number. Our
primitives separate them: the ability set $A$ is the opportunity, the chosen
bundle $(c,j)$ is the choice, and the profile $\mathbf y$ is the price
attached to opportunities. Holding the non-pecuniary content of a job fixed
in preferences while letting $\mathbf y$ vary lets the same task at different
pay---across regions, sectors, or time---be represented without altering how
the individual feels about the job itself.\footnote{%
The varying-$\mathbf y$ formulation is a modelling convenience, not a claim
that the classical model fixes pay. A sufficiently rich $\mathcal J$ with a
fixed profile could represent the same situations; we let the pecuniary
attribute vary while holding the non-pecuniary attribute fixed because it is
the more economical representation. Cf.\ the analogous point in
\citet{ManiquetMoramarco2026} that actual incomes do not deliver
preference-consistent comparisons across agents facing different prices.}

The discrete opportunity set is the categorical counterpart of the unequal
\emph{skills} of the fair-allocation literature
\cite{FleurbaeyManiquet1996,FleurbaeyManiquet2006}: where that literature
replaces a common budget line by an individual scalar skill $a_i\in\mathbb
R_+$, we replace it by an individual \emph{set} $A\subseteq\mathcal J$. The
difference is that a set of categorical jobs carries no order. There is no
labour line to be tangent to, and so the single-crossing structure on which
the continuous characterizations rely---and whose breakdown under
two-dimensional heterogeneity in skill \emph{and} preferences
\citet{FleurbaeyManiquet2006} themselves flag as the source of ranking
reversals and bunching---is simply unavailable. In its place, the laisser-faire
benchmark underlying $W^3$ and $W^{5,\bar A}$ is constructed by adding a
\emph{uniform subsidy} to pre-tax income across a set of jobs and asking which
job the individual would then choose. This is the discrete-opportunity
analogue of moving along a budget line, but it is an operation on a set, not a
point on a continuum.\footnote{%
We do not claim novelty for evaluating agents against opportunity sets as
such; the egalitarian-equivalence and equal-well-being-for-equal-budget
traditions do this. The novel primitive is the \emph{discrete, unordered,
agent-specific} feasible set $A$ together with a separately varying pay
profile $\mathbf y$.}

\subsection{Where the compensation--responsibility cut is drawn}
\label{sub:compresp}

Because opportunity and choice are fused in the classical wage, the
compensation--responsibility distinction there must be drawn along a single
dimension, which is the source of the recurring tension between compensating
low skill and respecting heterogeneous preferences over labour. Separating
the primitives relocates the cut. In our framework the question
``responsibility for what, compensation for what?'' becomes the question
``\emph{relative to whose opportunity set is well-being measured?}'' and the
two coherent answers are precisely our two measures.

\paragraph{The two measures are the two sides of the cut.}
The own-set measure $W^3$ evaluates each individual relative to her
\emph{own} ability set $A$: it holds her responsible for the opportunities she
faces and for her choice within them. The reference-set measure
$W^{5,\bar A}$ evaluates every individual relative to a \emph{common}
reference set $\bar A$: it compensates for differences in opportunity by
refusing to let the actual set $A$ affect the measured level. The axis that
distinguishes the two measures in our classification---own $A$ versus
reference $\bar A$---\emph{is} the compensation--responsibility axis, now
drawn over opportunity sets rather than over a wage scalar. This is made
visible by a single entry in the classification table: $W^{5,\bar A}$
satisfies Responsibility for Reference Abilities while $W^3$ does not, and
that one row is the formal expression of the distinction.

\paragraph{A structural by-product: responsibility absorbs the
irrelevant-jobs principle.}
The discrete model also yields a sharpening of the cut that the continuous
model cannot state. We show (\S\ref{sec:W3}, Proposition on the redundancy of
Independence of Irrelevant Jobs) that Independence of Irrelevant Jobs is
\emph{not} an independent normative commitment once one imposes Full
Responsibility together with the laisser-faire normalization: under the
maintained Representation axiom, the principle that feasible-but-dominated
jobs are normatively inert is a \emph{consequence} of holding individuals
responsible for their choices within a fixed opportunity set. In the
continuous model this statement has no analogue, because every labour level is
feasible and chosen by someone; there is no notion of a feasible-but-never-taken
alternative whose deletion must be shown not to matter. That Independence of
Irrelevant Jobs is derivable rather than postulated is, in this sense, a
result about the internal structure of responsibility in a discrete-jobs
world.\footnote{%
The independence of \emph{irrelevant commodities} in the consumption setting
\textcite{FleurbaeyTadenuma2007} is a structural cousin of this principle but
concerns commodities in a divisible bundle, not feasible-but-unchosen jobs in
a categorical set; we invoke it as an analogy, not an identity.}

\subsection{Feasibility phenomena with no continuous analogue}
\label{sub:feasibility}

In a continuum every labour level $\ell\in[0,1]$ is feasible, so a battery of
distinctions that are central to our analysis cannot even be posed there. A
job may be \emph{infeasible} (outside $A$), \emph{feasible but irrelevant}
(in $A$ yet dominated at every tax, hence never chosen), or simply
\emph{not currently held}. These are three different situations, and the
continuous model collapses all of them, because the only sense in which a
labour level is ``not held'' is that the individual chose a different
one.\footnote{%
The three-way distinction parallels, but is not imported from,
\citet{FleurbaeyTadenuma2007}, \S6, who separate goods that are unavailable,
available-but-unconsumed, and in shortage, and who note that how finely
alternatives are individuated is a borderline of model specification. We make
the analogous individuation choices for jobs explicitly.}
The axioms that govern these distinctions---Job Duplication Invariance, Job
Neutrality, Independence of Irrelevant Jobs, and Independence of Preferences
over Infeasible Jobs---are accordingly new to this setting and should be read
as such, not as restatements of axioms from the divisible-goods literature.

The staying-home option is the cleanest instance. We include in every ability
set a job $o$ representing non-employment, present in all $A$ by construction.
This has a precise precedent in the assignment literature: the null object of
\citet{Maniquet2008}, always available and conferring only money, is the
structural ancestor of our staying-home job, and it is what allows the measure
to be defined for individuals at the boundary of the labour market---the
unemployed and the non-participating---whom the continuous model can represent
only by the degenerate device of a zero wage above the current labour level
\textcite[cf.][]{FleurbaeyManiquet2007}.

\subsection{Implications for social evaluation and taxation}
\label{sub:swf}

A well-being measure is an input to social evaluation, not a social ordering
in itself. We therefore separate measurement from aggregation, following the
program of \citet{FleurbaeyManiquet2017} and \citet{fleurbaey_maniquet_2019}:
well-being is first measured agent by agent, and only then aggregated. This
separation is a methodological stance rather than a theorem of the present
model, and we note its one important qualification---that the choice of
aggregator is not fully independent of the choice of measure once fairness
properties are imposed on allocations
\textcite[][p.~785]{fleurbaey_maniquet_2019}, so that the responsibility content
of $W^3$ versus $W^{5,\bar A}$ does constrain which aggregators preserve the
intended fairness properties.

\paragraph{Why measure first, aggregate second.}
The separation is not merely tidy; it is a way around an impossibility. Direct
aggregation of ordinal, non-comparable preferences runs into dictatorship
results: even weak independence together with the Pareto principle forces the
social ranking to track a single individual unless the informational basis or
the domain is enriched \textcite{FleurbaeyTadenuma2007}. Constructing an
individual measure $W$ first---an object denominated in a common unit and
comparable across agents---and aggregating that, rather than aggregating raw
preferences, is precisely how the fair-allocation program sidesteps the
impossibility.

\paragraph{A natural aggregator, presented as such.}
When the individual measure is constructed from fairness axioms rather than
assumed, the social-ordering-function literature shows that the aggregator
those axioms permit is essentially the leximin criterion
\textcite{Maniquet2016,FleurbaeyManiquet2006,Valletta2009,ManiquetNeumann2021}.
We present maximin/leximin aggregation of our measures as the precedented and
axiomatically natural choice, \emph{not} as a result already established for
the discrete-jobs model; a full aggregation theorem in this setting is left to
future work, and the precedents themselves caution that it is the richness of
the preference domain, not the aggregation axioms alone, that forces extreme
inequality aversion---under quasi-linear preferences a broad class of concave
aggregators would qualify \textcite[][p.~213]{Maniquet2008}.

\paragraph{The policy direction is immediate in the measure's units.}
Because $W^3$ and $W^{5,\bar A}$ are denominated in the uniform subsidy $w$,
the measure is already in the unit of the policy instrument: a leximin
objective $\max\min_i W_i$ reads directly as ``raise the subsidy, or lower the
tax, of the worst-off type,'' which is the discrete-jobs analogue of the
one-dimensional read-off of a tax reform from the budget frontier in the
continuous model \textcite[][\S3.1]{Maniquet2016}. The characteristic shapes that
the continuous fair-tax model delivers---non-positive marginal rates at the
bottom, the largest subsidy at the lowest opportunity, a basic-income
rationale as the lowest opportunity degenerates
\textcite{FleurbaeyManiquet2006,FleurbaeyManiquet2007,ManiquetNeumann2021}---are
the natural conjectures for the discrete model and the right benchmark to
test against, but they are outputs of the continuous tax apparatus and would
have to be re-derived for the discrete-jobs optimal-tax problem; we do not
claim them as theorems here.

\subsection{Summary of the contribution}
\label{sub:summary}

The exercise is worth carrying out because separating opportunity, choice, and
price into the three primitives $A$, $(c,j)$, and $\mathbf y$ lets the model
state things the labour-line model cannot: that equal pay need not mean equal
well-being when jobs differ in content; that the compensation--responsibility
distinction is a choice of reference opportunity set, with $W^3$ and
$W^{5,\bar A}$ as its two coherent resolutions; that the irrelevant-jobs
principle is a consequence of responsibility rather than a separate axiom; and
that feasibility, infeasibility, and the staying-home margin are distinct
phenomena rather than a single labour quantity. None of these requires denying
that the continuous model is a powerful and well-understood benchmark; each
identifies a question that becomes expressible only once labour time is treated
as an attribute attached to a discrete, categorical, possibly infeasible job
rather than as a quantity chosen freely on a continuum.
